\documentclass[12pt]{book}
\usepackage{precalc1}

% --- solutions list environment: numbered to match exercises ---
\newcounter{soln}[section]
\newenvironment{solutions}{%
  \setcounter{soln}{0}%
  \begin{list}{\textbf{\arabic{soln}.}}{%
    \usecounter{soln}\setlength{\leftmargin}{2.2em}\setlength{\itemsep}{10pt}}%
}{\end{list}}

\begin{document}

\chapter*{Solutions Manual}
\markboth{Precalculus I — Solutions}{Precalculus I — Solutions}
\addcontentsline{toc}{chapter}{Solutions Manual}

\bigskip
\noindent\textit{Worked solutions to the end-of-chapter exercises, Weeks 1--10.
Conceptual answers give the reasoning expected; computational answers show the steps.
Sketching and lab exercises describe what a complete response contains.}

\section*{Week 1 — Functions}
\setcounter{section}{1}

\begin{solutions}
\item Domain $\{-1,0,2,5\}$, range $\{3,0,-3,1\}$. \emph{Not} a function: the input
$2$ is paired with both $3$ and $-3$, so one input determines two outputs.

\item (a) A circle fails; a vertical line through its interior meets it twice.
(b) A downward parabola passes; every vertical line meets it once.
(c) A sideways ``S'' fails; some vertical lines meet it three times.
(d) A single steep line passes. The test checks the defining property of a function:
that each input has exactly one output, since a vertical line collects all points
sharing one $x$-value.

\item $f(0)=3$ and $f(-2)=2(4)-(-2)+3=8+2+3=13$. For $f(t+1)$,
\[
f(t+1)=2(t+1)^2-(t+1)+3=2(t^2+2t+1)-t-1+3=2t^2+3t+4.
\]

\item For example $\{(1,2),(1,5)\}$ is not a function, since $1$ has two outputs.
Changing $(1,5)$ to $(2,5)$ gives $\{(1,2),(2,5)\}$, a function: the repeated input is
removed, so every input now has a single output.

\item Shifting left $4$ uses $h=-4$, stretching vertically by $3$ uses $A=3$, and
shifting down $1$ uses $k=-1$:
\[
g(x)=3(x+4)^2-1.
\]
The vertex of $x^2$ at the origin maps to $(h,k)=(-4,-1)$.

\item Read $g(x)=-\tfrac12\,|x+2|+5$ as $A=-\tfrac12$, $h=-2$, $k=5$. In order:
shift left $2$; compress vertically by $\tfrac12$ and reflect across the $x$-axis
(because $A<0$); shift up $5$. The corner of $|x|$ at the origin lands at
$(h,k)=(-2,5)$.

\item Track the point $(9,3)$ through $g(x)=2\sqrt{x-1}+4$. The input transforms by
the inside: $x:9\to 9-1=8$? No, undo carefully. The inside is $x-1$, so the parent's
input $9$ corresponds to $g$'s input $9+1=10$. The output transforms by the outside:
$y:3\to 2(3)+4=10$. So $(9,3)\mapsto(10,10)$. Confirm directly:
$g(10)=2\sqrt{10-1}+4=2(3)+4=10$. \checkmark

\item Constants outside $f$ act on the output and move the graph vertically as
expected, so $f(x)+3$ raises every output by $3$. Constants inside $f$ act on the
input and move the graph horizontally in reverse: $f(x-3)$ reaches the parent's value
at input $0$ when $x-3=0$, that is at $x=3$, so the graph shifts right. The two
``$+3$'' moves differ because one adjusts the output (after the function acts) and the
other adjusts the input (before the function acts).

\item The cloud is the data itself: scattered points with no exact rule. A linear fit
is a \emph{model}, a clean predictable line laid through the cloud to follow its
trend. The two are not equal; the model summarizes and forecasts the data without
passing through every point. ``Model'' names the chosen rule that stands in for the
data so one can compute with it.

\item With $A<0$ the graph reflects across the $x$-axis (the output sign flips), and
with $b<0$ it reflects across the $y$-axis (the input sign flips). The two reflections
are independent and involve the $x$-axis and the $y$-axis respectively.

\item A bowl opening upward with lowest point near $(2,-3)$ suggests the parent
$f(x)=x^2$. Start with $A=1$ (upward, unstretched), $b=1$, $h=2$ (vertex $x$), and
$k=-3$ (vertex $y$). The vertex of the model then sits at $(2,-3)$, matching the
cloud's lowest point; $A$ is tuned afterward to match how steeply the cloud rises.

\item \emph{Lab exercise.} Answers vary with the generated cloud. A complete response
records the parent family, the four fitted constants, the fit score, and one sentence
identifying the hardest constant (commonly $b$, since horizontal scaling is the least
visually obvious of the four).
\end{solutions}

\section*{Week 2 — Lines and Parabolas}
\setcounter{section}{2}

\begin{solutions}
\item Any two points on a line determine a right triangle with horizontal leg (run)
and vertical leg (rise). Because the line meets every horizontal at one fixed angle,
all such triangles have the same angles and are therefore similar. Similar triangles
have proportional sides, so the ratio rise$/$run is identical for every pair of
points; that common ratio is the slope.

\item Slope: $m=\dfrac{-6-2}{3-(-1)}=\dfrac{-8}{4}=-2$. Point-slope through $(-1,2)$:
$y-2=-2(x+1)$. Slope-intercept: $y=-2x-2+2=-2x$, so $f(x)=-2x$. Check: $f(-1)=2$ and
$f(3)=-6$. \checkmark

\item Point-slope: $y-1=-\tfrac23(x-4)$. Slope-intercept:
$y=-\tfrac23 x+\tfrac83+1=-\tfrac23 x+\tfrac{11}{3}$. The $y$-intercept is
$\tfrac{11}{3}$. The $x$-intercept solves $0=-\tfrac23 x+\tfrac{11}{3}$, giving
$x=\tfrac{11}{2}$.

\item (a) Through $(2,5)$ and $(2,-1)$: the run is $0$, so the slope is undefined; this
is the vertical line $x=2$, which is \emph{not} a function. (b) Through $(0,4)$ and
$(7,4)$: slope $0$, a function (the horizontal line $y=4$). (c) Through $(1,1)$ and
$(3,5)$: slope $\dfrac{5-1}{3-1}=2$, a function.

\item A parabola is the set of points equidistant from a fixed point (the focus) and a
fixed line (the directrix). The vertex is the point of the parabola on the axis; its
distance to the focus equals its distance to the directrix by definition. Since both
distances are measured along the axis, the vertex sits exactly halfway between focus
and directrix.

\item Rays arriving parallel to the axis reflect off the curve and all pass through the
focus. Two devices: a satellite dish or solar collector, where incoming parallel rays
are concentrated \emph{at} the focus; and a headlight or flashlight reflector, where
light from a source \emph{at} the focus is sent out as a parallel beam. The directions
are reversed in the two cases.

\item Substitute each point into $f(x)=ax^2+bx+c$. From $(0,1)$: $c=1$. From
$(-1,6)$: $a-b+c=6$, so $a-b=5$. From $(2,3)$: $4a+2b+c=3$, so $4a+2b=2$, i.e.
$2a+b=1$. Adding $a-b=5$ and $2a+b=1$ gives $3a=6$, so $a=2$, then $b=-3$. Thus
$f(x)=2x^2-3x+1$. Check $(-1)$: $2+3+1=6$. \checkmark

\item Factor $3$ from the variable terms: $3(x^2+4x)+7$. Half of $4$, squared, is $4$:
$3(x^2+4x+4-4)+7=3(x+2)^2-12+7=3(x+2)^2-5$. The vertex is $(-2,-5)$. By formula,
$h=-\dfrac{b}{2a}=-\dfrac{12}{6}=-2$, agreeing.

\item Factor: $x^2-2x-15=(x-5)(x+3)$, roots $5$ and $-3$. Axis of symmetry:
$x=\dfrac{5+(-3)}{2}=1$. Vertex $y$-value: $f(1)=1-2-15=-16$, so the vertex is
$(1,-16)$. The $y$-intercept is $f(0)=-15$. A correct sketch is an upward parabola
crossing at $-3$ and $5$, with vertex $(1,-16)$ and $y$-intercept $(0,-15)$.

\item Discriminant $b^2-4ac$. (a) $x^2-4x+4$: $16-16=0$, one (doubled) real root at
$x=2$. (b) $2x^2+x+3$: $1-24=-23<0$, no real roots. (c) $x^2-5x+2$: $25-8=17>0$, two
real roots, $x=\dfrac{5\pm\sqrt{17}}{2}\approx 4.56$ and $0.44$.

\item With roots $-4$ and $2$, $f(x)=a(x+4)(x-2)$. Then $f(0)=a(4)(-2)=-8a=-16$, so
$a=2$. The three forms:
\[
\text{factored } 2(x+4)(x-2),\quad
\text{standard } 2x^2+4x-16,\quad
\text{vertex } 2(x+1)^2-18,
\]
the last from $h=\dfrac{-4+2}{2}=-1$ and $k=f(-1)=2-4-16=-18$.

\item The vertex sits at $h=-\dfrac{b}{2a}$. The root-sum relation gives
$r_1+r_2=-\dfrac{b}{a}$, so the average of the roots is
$\dfrac{r_1+r_2}{2}=-\dfrac{b}{2a}=h$. Thus the vertex's $x$-coordinate is exactly the
average of the roots. Geometrically, a parabola is symmetric about the vertical line
through its vertex, and the two roots lie at equal distances on either side of that
axis, so their midpoint is the axis itself.
\end{solutions}

\section*{Week 3 — Polynomials and Their Shape}
\setcounter{section}{3}

\begin{solutions}
\item End behavior from the leading term. (a) $3x^4$: even degree, positive lead, both
ends up. (b) $-x^3$: odd, negative lead, up on the left and down on the right.
(c) $-4x^6$: even, negative lead, both ends down. (d) $x^5$: odd, positive lead, down
on the left and up on the right.

\item For large $|x|$, a higher power outgrows every lower power: $x^n$ eventually
dwarfs $x^{n-1}, \dots, x, 1$. So the leading term decides the ends, and the constant
and middle coefficients, attached to lower powers, become negligible there. They shift
and bend the middle of the graph but cannot change where the ends go.

\item Roots: $-1$ (multiplicity $1$), $2$ (multiplicity $3$), $4$ (multiplicity $1$).
Degree $1+3+1=5$. The graph crosses at all three, since every multiplicity is odd; at
$x=2$ the crossing flattens against the axis because the multiplicity is $3$.

\item Just left of $r$ the factor $(x-r)$ is negative; just right it is positive. If
$m$ is odd, $(x-r)^m$ keeps the sign of $(x-r)$ (a negative raised to an odd power
stays negative), so it flips sign across $r$. If $m$ is even, $(x-r)^m\ge0$ on both
sides (an even power is nonnegative), so the sign does not change.

\item Roots $-3,1,2$ split the line into four intervals.
\[
\begin{array}{c|cccc}
\text{interval} & (-\infty,-3) & (-3,1) & (1,2) & (2,\infty)\\\hline
\text{test }x & -4 & 0 & 1.5 & 3\\
x+3 & - & + & + & +\\
x-1 & - & - & + & +\\
x-2 & - & - & - & +\\\hline
f & - & + & - & +
\end{array}
\]
The sign changes at each root (all simple).

\item For $f(x)=-(x+1)^2(x-3)$ the roots are $-1$ (multiplicity $2$) and $3$
(multiplicity $1$). Intervals: $f$ is $+$ on $(-\infty,-1)$, $+$ on $(-1,3)$, and $-$
on $(3,\infty)$. The sign does \emph{not} change at $x=-1$ (even multiplicity: a
touch) and \emph{does} change at $x=3$ (odd multiplicity: a crossing). The overall
leading sign is negative because of the $-1$ in front.

\item $f(x)=(x-1)(x+2)^2$. Degree $3$, positive lead: down on the left, up on the
right. Roots: cross at $1$ (simple), touch at $-2$ (double). Sign chart: $-$ on
$(-\infty,-2)$, $-$ on $(-2,1)$, $+$ on $(1,\infty)$. The curve rises from the lower
left, touches the axis at $x=-2$ and turns back down, continues below to cross up at
$x=1$, then climbs to the upper right.

\item Degree-five, positive lead: down-left to up-right. Roots: simple at $-2$ and $3$
(crossings), triple at $0$ (a flattened crossing). At the triple root the graph
crosses the axis but flattens noticeably against it, the visual signature of
multiplicity three, before continuing through.

\item Even degree with simple roots at $-1,4$ (multiplicity $1$ each) and a double root
at $2$ gives multiplicities summing to at least $1+1+2=4$, so the lowest possible
degree is $4$ (even, consistent). With negative lead, both ends point down. The curve
comes up from the lower left, crosses at $-1$, rises, touches at $2$ and turns, crosses
at $4$, and falls to the lower right.

\item The sign changes at $-1$ (from $+$ to $-$) and at $2$ (from $-$ to $+$), so both
roots have \emph{odd} multiplicity. The lowest-degree factored form using simple roots
is $f(x)=(x+1)(x-2)$ up to a positive constant; its sign pattern is $+,-,+$, matching.

\item Without calculus, the sign chart and roots fix which side of the axis the curve
occupies and where it crosses or touches, but not the exact height of the turning
point between the roots. Both students' graphs can show the correct shape (touch at
$x=1$, cross at $x=4$, correct ends) while differing in that height. Calculus, through
the derivative, is the tool that pins the turning-point height exactly.

\item \emph{Lab exercise.} Answers vary. A complete response records the factored form,
the fit score, and one sentence noting that switching a multiplicity from even to odd
turns a touch into a crossing at that root (or vice versa).
\end{solutions}

\section*{Week 4 — Rational Functions}
\setcounter{section}{4}

\begin{solutions}
\item Zeros at $x=4$ and $x=-2$ (numerator roots); vertical asymptotes at $x=1$ and
$x=-5$ (denominator roots). No factor is shared between numerator and denominator, so
there are no holes and no cancellation.

\item The factor $(x+2)$ is common, so it cancels, leaving $\dfrac{x-3}{x-5}$ with a
hole at $x=-2$. The zero is $x=3$ and the vertical asymptote is $x=5$.

\item Near a numerator root the fraction is (small)$/$(ordinary), which is small, so the
graph sits near the $x$-axis: a zero. Near a denominator root the fraction is
(ordinary)$/$(small), which is enormous, so the graph races off toward $\pm\infty$
along a vertical line: an asymptote.

\item Horizontal asymptotes. (a) Equal degrees: $y=\tfrac{2}{1}=2$. (b) Equal degrees:
$y=\tfrac{5}{1}=5$. (c) Numerator degree below denominator: $y=0$. (d) Numerator degree
above denominator: no horizontal asymptote.

\item $\dfrac{6x^2+x}{3x^2-12}$. Equal degrees, so $y=\tfrac{6}{3}=2$. Zeros from
$x(6x+1)=0$: $x=0$ and $x=-\tfrac16$. Vertical asymptotes from $3(x^2-4)=0$: $x=2$ and
$x=-2$.

\item Dividing $x^2-5x+6$ by $x-1$: quotient $x-4$, remainder $2$, so
$f(x)=x-4+\dfrac{2}{x-1}$. The slant asymptote is $y=x-4$.

\item Dividing $2x^2+3x-1$ by $x+2$: quotient $2x-1$, remainder $1$, so the slant
asymptote is $y=2x-1$. The vertical asymptote is $x=-2$. The zeros solve
$2x^2+3x-1=0$: $x=\dfrac{-3\pm\sqrt{17}}{4}\approx 0.28$ and $-1.78$.

\item When the numerator degree exceeds the denominator degree, the ratio of leading
terms grows without bound as $|x|\to\infty$, so the values cannot level off at a finite
height; there is no horizontal asymptote. The end-behavior curve is the polynomial
quotient from long division, and it is a straight line exactly when the numerator
degree is one more than the denominator degree.

\item $\dfrac{x-1}{(x+2)(x-3)}$. Features $-2,1,3$.
\[
\begin{array}{c|cccc}
\text{interval} & (-\infty,-2) & (-2,1) & (1,3) & (3,\infty)\\\hline
\text{test }x & -3 & 0 & 2 & 4\\
x-1 & - & - & + & +\\
x+2 & - & + & + & +\\
x-3 & - & - & - & +\\\hline
f & - & + & - & +
\end{array}
\]

\item $\dfrac{x^2-4}{x^2-1}=\dfrac{(x-2)(x+2)}{(x-1)(x+1)}$; no common factors, so no
holes. Zeros at $x=\pm2$ (both crossings). Vertical asymptotes at $x=\pm1$. Equal
degrees with leading coefficients $1$, so horizontal asymptote $y=1$. The sign chart
over $-2,-1,1,2$ places the curve above the axis on the outer intervals (approaching
$y=1$) and produces the escapes beside $x=\pm1$. The sketch is symmetric: outer
branches approach $y=1$ from above, the central branch passes through $\pm2$ with a
small peak at $(0,4)$? No: $f(0)=\tfrac{-4}{-1}=4$, so the central piece peaks near
$(0,4)$ — adjust: the central region lies between the asymptotes and the curve dips
below the axis between the zeros. A correct sketch shows the two outer branches and the
behavior near each asymptote, with $f(0)=4$.

\item Zeros at $0$ and $3$ need numerator factors $x(x-3)$; vertical asymptotes at
$-1$ and $2$ need denominator factors $(x+1)(x-2)$. The horizontal asymptote $y=1$
requires equal degrees with equal leading coefficients, which $\dfrac{x(x-3)}{(x+1)(x-2)}$
satisfies (both degree $2$, leading coefficient $1$). So
$f(x)=\dfrac{x(x-3)}{(x+1)(x-2)}$.

\item \emph{Lab exercise.} Answers vary. A complete response records the factored form,
the degree ratio, the fit score, and one sentence noting that raising the numerator
degree past the denominator degree replaces the horizontal asymptote with a slant or
polynomial end-behavior curve.
\end{solutions}

\section*{Week 5 — Inverses and Radicals}
\setcounter{section}{5}

\begin{solutions}
\item Domains. (a) $\dfrac{x}{x^2-9}$: all reals except $x=\pm3$. (b) $\sqrt{2x-6}$:
need $2x-6\ge0$, so $x\ge3$. (c) $\sqrt[3]{x+1}$: all reals (odd root). (d)
$\dfrac{1}{\sqrt{x-4}}$: need $x-4>0$ (strictly, to avoid both the negative radicand
and a zero denominator), so $x>4$.

\item Ranges. (a) $x^2-4$: a square is $\ge0$, so the range is $y\ge-4$. (b)
$\sqrt{x}+2$: $\sqrt{x}\ge0$, so the range is $y\ge2$. (c) $-|x|$: $|x|\ge0$, so
$-|x|\le0$, range $y\le0$.

\item One-to-one by the horizontal line test. (a) $3x+2$: yes (a non-horizontal line).
(b) $x^2-1$: no (a horizontal line meets it twice). (c) $x^3$: yes. (d) $|x|$: no (the
V is met twice by horizontals above the vertex).

\item An inverse must send each output back to the single input that produced it. If
the function is not one-to-one, some output $y$ comes from two different inputs, and
the inverse would have to send $y$ to both — impossible for a function. The shared
output is exactly the obstruction.

\item Inverses. (a) $f(x)=3x-7$: $f^{-1}(x)=\dfrac{x+7}{3}$, domain all reals. (b)
$f(x)=\dfrac{x+2}{5}$: $f^{-1}(x)=5x-2$, domain all reals. (c) $f(x)=x^3+1$:
$f^{-1}(x)=\sqrt[3]{x-1}$, domain all reals.

\item Restricting $f(x)=x^2$ to $x\ge0$ makes it one-to-one, with inverse
$f^{-1}(x)=\sqrt{x}$. Check: $f^{-1}(f(x))=\sqrt{x^2}=|x|=x$ for $x\ge0$. The inverse
has domain $x\ge0$ and range $y\ge0$ (the original domain).

\item $\sqrt[3]{x}$ inverts $x^3$, which takes every real value including negatives
(e.g. $(-2)^3=-8$), so its inverse is defined for all reals. $\sqrt{x}$ inverts $x^2$,
which is never negative, so there is no real input whose square is negative; the domain
is restricted to $x\ge0$. The difference is the parity of the power being inverted.

\item Sketch: $f(x)=x^2$ on $x\ge0$ is the right half of the upward parabola, starting
at the origin; $f^{-1}(x)=\sqrt{x}$ is the upper half of a sideways parabola, also from
the origin. Each is the mirror image of the other across the dashed line $y=x$, since
the point $(a,a^2)$ on $f$ corresponds to $(a^2,a)$ on $f^{-1}$.

\item (a) $\sqrt{2x+3}=x$. Square: $2x+3=x^2$, so $x^2-2x-3=0$, $(x-3)(x+1)=0$,
candidates $3,-1$. Check: $x=3$ gives $\sqrt9=3$ \checkmark; $x=-1$ gives $\sqrt1=1\ne-1$,
rejected. Solution $x=3$. (b) $\sqrt{x+7}=x+1$. Square: $x+7=x^2+2x+1$, so
$x^2+x-6=0$, $(x+3)(x-2)=0$, candidates $-3,2$. Check: $x=2$ gives $\sqrt9=3=2+1$
\checkmark; $x=-3$ gives $\sqrt4=2\ne-2$, rejected. Solution $x=2$.

\item Reflecting across $y=x$ sends a point $(a,b)$ to $(b,a)$. A point on both $f$ and
$f^{-1}$ is fixed by this reflection, which happens only when $a=b$, i.e. on the line
$y=x$. So any intersection of a function with its inverse lies on that line.

\item $f(x)=\sqrt{x-2}+1$: domain $x\ge2$ (radicand $\ge0$), range $y\ge1$ (the
$\sqrt{\,}$ is $\ge0$, shifted up $1$). Inverse: set $y=\sqrt{x-2}+1$, swap, solve:
$x-1=\sqrt{y-2}$, so $y=(x-1)^2+2$, giving $f^{-1}(x)=(x-1)^2+2$ with domain $x\ge1$.
The inverse's domain $x\ge1$ equals the range of $f$, as required.

\item $f(x)=\dfrac1x$. Swap and solve: $x=\dfrac1y$ gives $y=\dfrac1x$, so
$f^{-1}(x)=\dfrac1x=f(x)$. The function is its own inverse. Its graph is therefore
symmetric across the line $y=x$, since reflecting it there returns the same curve.
\end{solutions}

\section*{Week 6 — Exponential Functions}
\setcounter{section}{6}

\begin{solutions}
\item (a) $5,10,20,40$: constant ratio $2$, exponential $f(x)=5\cdot2^{x}$. (b)
$5,8,11,14$: constant difference $3$, linear $f(x)=3x+5$. (c) $100,20,4$: constant
ratio $\tfrac15$, exponential $f(x)=100\left(\tfrac15\right)^{x}$.

\item Both $3^x$ and $3^{-x}$ pass through $(0,1)$. The graph of $3^x$ rises to the
right (growth); $3^{-x}$ is its mirror across the $y$-axis, falling to the right
(decay). Each has the $x$-axis ($y=0$) as a horizontal asymptote.

\item (a) $5^{-x}=\left(\tfrac15\right)^{x}$, decay. (b) $\left(\tfrac14\right)^{x}$ is
already positive-base positive-exponent decay. (c) $2^{-3x}=\left(\tfrac18\right)^{x}$,
since $2^{-3}=\tfrac18$, decay.

\item Each finer compounding adds a little more growth, but each added amount is
smaller than the last, so the running total is trapped beneath a fixed ceiling rather
than increasing without bound. It approaches $e\approx2.71828$.

\item $A(10)=1000\,e^{0.04\cdot10}=1000\,e^{0.4}\approx 1000(1.49182)\approx
\$1491.82$.

\item $4^{x}=e^{(\ln 4)x}$ with $\ln 4\approx1.3863$, so $f(x)\approx e^{1.3863\,x}$.

\item $0.8^{t}=e^{(\ln 0.8)t}$ with $\ln 0.8\approx-0.2231$. The continuous rate is
negative, confirming decay: the quantity shrinks at about $22.3\%$ continuous rate per
unit time.

\item Each $10$ years is one half-life: $64\to32$ ($t=10$), $32\to16$ ($t=20$),
$16\to8$ ($t=30$). Remaining masses: $32$, $16$, $8$ grams.

\item $M(t)=M_0\left(\tfrac12\right)^{t/1600}$. After $4800$ years,
$t/T=4800/1600=3$, so $\left(\tfrac12\right)^{3}=\tfrac18$ remains.

\item $\left(\tfrac12\right)^{n}=\tfrac18$ gives $n=3$ half-lives. With $T=30$ years,
the age is $3\times30=90$ years.

\item A radioactive element loses a fixed \emph{fraction} of its mass in each fixed
time span (half in each half-life), which is constant-ratio behavior, the defining
property of an exponential function. Subtracting a fixed amount each year would be
linear and would reach zero in finite time, which decay never does.

\item After $8$ years: the $2$-year isotope has passed $8/2=4$ half-lives, leaving
$\left(\tfrac12\right)^{4}=\tfrac1{16}$; the $4$-year isotope has passed $8/4=2$
half-lives, leaving $\left(\tfrac12\right)^{2}=\tfrac14$. The shorter-half-life isotope
has decayed further, because it halves more often in the same time.
\end{solutions}

\section*{Week 7 — Logarithms}
\setcounter{section}{7}

\begin{solutions}
\item (a) $\log_3 81=4$ ($3^4=81$). (b) $\log_{10}0.01=-2$ ($10^{-2}=0.01$). (c)
$\log_2\tfrac1{16}=-4$ ($2^{-4}=\tfrac1{16}$). (d) $\log_7 1=0$. (e) $\ln e^4=4$.

\item (a) $5^3=125 \iff \log_5 125=3$. (b) $\log_2 32=5 \iff 2^5=32$. (c)
$e^0=1 \iff \ln 1=0$.

\item The graphs of $2^x$ and $\log_2 x$ are reflections across $y=x$. The exponential
passes through $(0,1)$ with horizontal asymptote $y=0$; the logarithm passes through
$(1,0)$ with vertical asymptote $x=0$ (the $y$-axis).

\item (a) $\log_b(x^2yz^3)=2\log_b x+\log_b y+3\log_b z$. (b)
$\log\dfrac{\sqrt{x}}{y^4}=\tfrac12\log x-4\log y$. (c)
$\ln\dfrac{e^2 x}{y}=2+\ln x-\ln y$ (since $\ln e^2=2$).

\item (a) $\log x+3\log y=\log(xy^3)$. (b) $2\ln x-\tfrac12\ln y=\ln\dfrac{x^2}{\sqrt
y}$. (c) $\log_b 6-\log_b 2=\log_b 3$.

\item The laws apply to products, quotients, and powers, never to sums; $\log_b(M+N)$
has no expansion. The expression $\log_b M+\log_b N$ equals $\log_b(MN)$, the logarithm
of the product, which is a different quantity from $\log_b(M+N)$.

\item By change of base: (a) $\log_2 50=\dfrac{\ln 50}{\ln 2}\approx5.6439$. (b)
$\log_5 100=\dfrac{\ln 100}{\ln 5}\approx2.8614$. (c)
$\log_3 7=\dfrac{\ln 7}{\ln 3}\approx1.7712$.

\item (a) $2^x=15$: $x=\dfrac{\ln 15}{\ln 2}\approx3.9069$. (b) $5^x=2$:
$x=\dfrac{\ln 2}{\ln 5}\approx0.4307$. (c) $e^x=10$: $x=\ln 10\approx2.3026$.

\item (a) $\log_2 16=4$ bits. (b) $\log_2 1000=\dfrac{\ln 1000}{\ln 2}\approx9.97$
bits. (c) $\log_2 6\approx2.58$ bits.

\item $p=\tfrac1{64}$: information $-\log_2\tfrac1{64}=\log_2 64=6$ bits. A rarer event
has smaller $p$, hence larger $\log_2\tfrac1p$, so its occurrence resolves more
uncertainty and carries more information.

\item A password chosen from $2^{40}$ equally likely strings carries
$\log_2 2^{40}=40$ bits. By the product law, the bits of independent characters add:
$\log_2(2^k\cdot2^m)=\log_2 2^k+\log_2 2^m=k+m$, so each character contributes its own
bits to the total.

\item Total outcomes $6\times4=24$, carrying $\log_2 24\approx4.585$ bits. The parts
add: $\log_2 6+\log_2 4\approx2.585+2=4.585$, matching, because the product law of
logarithms is the additivity of information across independent sources.
\end{solutions}

\section*{Week 8 — Implicit Curves and Conics}
\setcounter{section}{8}

\begin{solutions}
\item On $x^2+y^2=25$? (a) $(0,5)$: $0+25=25$ \checkmark. (b) $(3,4)$: $9+16=25$
\checkmark. (c) $(-4,3)$: $16+9=25$ \checkmark. (d) $(2,2)$: $4+4=8\ne25$, not on the
circle.

\item Solving $x^2+y^2=9$ gives $y=\sqrt{9-x^2}$ (upper semicircle) and
$y=-\sqrt{9-x^2}$ (lower semicircle), each with domain $-3\le x\le3$. Together they
recombine to the full circle, since every point of the circle has $y$ either
nonnegative (upper) or nonpositive (lower).

\item For $y^2=x^3-x$: at $x=-1,0,1$ the right side is $0$, so $y=0$, giving the points
$(-1,0)$, $(0,0)$, $(1,0)$. At $x=\tfrac14$, $x^3-x=\tfrac1{64}-\tfrac14<0$, so $y^2<0$
is impossible: no real point there.

\item Wherever $x^3-x>0$, the equation gives $y=\pm\sqrt{x^3-x}$, two opposite values,
so a vertical line meets the curve twice and it fails the vertical line test. The real
locus has two pieces: a bounded oval over $-1\le x\le0$ and an unbounded branch over
$x\ge1$, with a gap on $(0,1)$ where $x^3-x<0$.

\item At $x=3$: $y^2=27-3=24$, so $y=\pm\sqrt{24}=\pm2\sqrt6$. The points are
$(3,2\sqrt6)$ and $(3,-2\sqrt6)$. Check: $(2\sqrt6)^2=24=3^3-3$ \checkmark.

\item The radius is the distance from the origin to $(6,8)$:
$\sqrt{6^2+8^2}=\sqrt{100}=10$. The circle is $x^2+y^2=100$.

\item $\dfrac{x^2}{25}+\dfrac{y^2}{16}=1$: $a^2=25$, $b^2=16$, so $a=5$, $b=4$.
Intercepts $(\pm5,0)$ and $(0,\pm4)$. The sketch is an ellipse stretching to $\pm5$
horizontally and $\pm4$ vertically.

\item $\dfrac{x^2}{9}-\dfrac{y^2}{4}=1$: $a=3$, $b=2$. Vertices $(\pm3,0)$; asymptotes
$y=\pm\dfrac{b}{a}x=\pm\dfrac23 x$. The two branches open left and right, hugging the
asymptotes.

\item Classify. (a) $3x^2+3y^2=12$: equal coefficients, circle. (b) $4x^2+y^2=16$:
unequal positive, ellipse. (c) $y^2-x^2=1$: opposite signs, hyperbola. (d)
$x-2y^2=0$: one squared term, parabola.

\item For large $x$, the constant $1$ is negligible beside the large squared terms, so
$\dfrac{x^2}{a^2}-\dfrac{y^2}{b^2}=1$ approaches $\dfrac{x^2}{a^2}=\dfrac{y^2}{b^2}$,
which rearranges to $y=\pm\dfrac{b}{a}x$. The branches approach these asymptote lines
ever more closely as $x$ grows.

\item Dividing $x$ by $a$ and $y$ by $b$ before squaring replaces $x^2+y^2=1$ with
$\dfrac{x^2}{a^2}+\dfrac{y^2}{b^2}=1$. This stretches the unit circle by a factor $a$
horizontally and $b$ vertically: the curve now reaches $\pm a$ along the $x$-axis and
$\pm b$ along the $y$-axis, so $a$ and $b$ are the semi-axes.

\item $4x^2-y^2=1$, i.e. $\dfrac{x^2}{1/4}-\dfrac{y^2}{1}=1$. Opposite signs: a
hyperbola. Here $a^2=\tfrac14$ so $a=\tfrac12$, vertices $\left(\pm\tfrac12,0\right)$;
$b=1$, asymptotes $y=\pm\dfrac{b}{a}x=\pm2x$.
\end{solutions}

\section*{Week 9 — Systems of Linear Equations}
\setcounter{section}{9}

\begin{solutions}
\item (a) Two lines with different slopes cross at one point: a unique solution. (b)
Two distinct lines with equal slope are parallel and never meet: no solution. (c) One
line written two ways is a single line: every point solves both, infinitely many
solutions.

\item From the first, $y=5-x$. Substitute: $2x-(5-x)=1$, so $3x=6$, $x=2$, then $y=3$.
Check: $2+3=5$ and $2(2)-3=1$ \checkmark. Solution $(2,3)$.

\item From the first, $y=3x-2$. Substitute into $6x-2y=10$: $6x-2(3x-2)=10$ gives
$4=10$, a contradiction. No solution; the lines are parallel.

\item The second equation is $-2$ times the first ($-2x+4y=-8$ is $x-2y=4$ scaled), so
it adds nothing. Substituting $x=4+2y$ into the second yields $-8=-8$, always true.
Infinitely many solutions: with $y=t$, the set is $(4+2t,\,t)$.

\item A linear system's solution set is the intersection of lines or planes, which are
flat objects. Flat objects intersect in a flat object: a point, the empty set, or an
entire line or plane. "Exactly two points" is not a flat set, so it cannot be a
solution set.

\item Two arrangements giving no solution: (i) two of the three planes are parallel and
distinct, so they share no point; (ii) the three planes form a triangular prism, each
pair meeting in a line but with the three lines parallel and distinct, so no point lies
on all three.

\item Back-substitute. $4z=8$ gives $z=2$. Then $y-z=1$ gives $y=3$. Then
$x-y+2z=5$ gives $x=5+3-4=4$. Solution $(4,3,2)$.

\item Eliminate. $R_2-2R_1$: $-y-3z=-5$. $R_3-R_1$: $-3y+2z=7$. From the first of
these, $y=5-3z$; substitute: $-3(5-3z)+2z=7$, so $-15+11z=7$, $z=2$, then $y=-1$, then
$x=4-(-1)-2=3$. Solution $(3,-1,2)$. Check: $3-1+2=4$ \checkmark.

\item The second equation is twice the first, so it is redundant. The system reduces to
$x+y+z=2$ and $x-y+z=0$. Subtracting, $2y=2$, so $y=1$; then $x+z=1$. With $z=t$ free,
$x=1-t$, giving solutions $(1-t,\,1,\,t)$: infinitely many, a line where the two
distinct planes meet.

\item Let $a,c$ be adult and child tickets. $a+c=200$ and $12a+8c=2000$. From the
first, $c=200-a$; substitute: $12a+8(200-a)=2000$, so $4a+1600=2000$, $a=100$, then
$c=100$. One hundred of each.

\item Let $x,y,z$ be pounds at \$8, \$10, \$14. Total $x+y+z=10$; value
$8x+10y+14z=11(10)=110$; equal cheaper beans $x=y$. Substituting $x=y$:
$2x+z=10$ and $18x+14z=110$. From the first $z=10-2x$; substitute:
$18x+14(10-2x)=110$, so $-10x+140=110$, $x=3$, then $y=3$, $z=4$. Three pounds each of
the \$8 and \$10 beans and four of the \$14.

\item Adding a multiple of one equation to another produces a new equation that holds
exactly when both originals do, and the step is reversible (subtract the same
multiple), so the solution set is unchanged. Scaling by zero is forbidden because it
replaces an equation with $0=0$, discarding its information and potentially enlarging
the solution set.
\end{solutions}

\section*{Week 10 — Matrices and Gaussian Elimination}
\setcounter{section}{10}

\begin{solutions}
\item
\[
\left[\begin{array}{ccc|c}
3&-1&2&5\\ 1&4&-1&0\\ 2&1&1&7
\end{array}\right]
\]
The first column holds the $x$-coefficients, the second the $y$-coefficients, the
third the $z$-coefficients, and the column after the bar the constants.

\item (i) \emph{Swap} two rows — reorder the equations. (ii) \emph{Scale} a row by a
nonzero constant — multiply an equation through by that constant. (iii) \emph{Replace}
a row by itself plus a multiple of another row — add a multiple of one equation to
another. Each leaves the solution set unchanged.

\item Replace $R_2$ by $R_2-2R_1$:
\[
\left[\begin{array}{cc|c}1&3&5\\ 2&1&4\end{array}\right]
\rightarrow
\left[\begin{array}{cc|c}1&3&5\\ 0&-5&-6\end{array}\right].
\]

\item $R_2-3R_1$ gives $\left[\begin{array}{cc|c}1&2&5\\ 0&-5&-10\end{array}\right]$.
The bottom row reads $-5y=-10$, so $y=2$; then $x+2(2)=5$, $x=1$. Solution $(1,2)$.

\item $R_2-2R_1$: $-y-3z=-5$. $R_3-R_1$: $-3y+2z=7$. Eliminating $y$ (e.g. $R_3-3R_2$
on these) gives $11z=22$, so $z=2$; then $y=-1$ and $x=3$. Solution $(3,-1,2)$. Check:
$3-1+2=4$ \checkmark.

\item $R_2+3R_1$ gives $\left[\begin{array}{cc|c}1&-2&3\\ 0&0&14\end{array}\right]$.
The bottom row reads $0=14$, a contradiction, so there is no solution; the row of zero
coefficients with nonzero constant reveals it.

\item $R_2-\tfrac32 R_1$ (or $2R_2-3R_1$) gives a bottom row of all zeros,
$0=0$. The variable $y$ is free; from $2x+4y=6$, $x=3-2y$. With $y=t$, the solution set
is $(3-2t,\,t)$.

\item A leading entry in every variable column pins down each variable: unique
solution. A row $[\,0\ \cdots\ 0\mid c\,]$ with $c\neq0$ reads $0=c$, a contradiction:
no solution. A row of all zeros leaves a variable without a leading entry, hence free:
infinitely many.

\item The row $[\,0\ 0\ 0\mid 5\,]$ reads $0=5$, which is false, so the system has no
solution.

\item Atom counts: Fe gives $a=2c$, O gives $2b=3c$. Take $c=t$: $a=2t$,
$b=\tfrac{3t}{2}$. The smallest whole numbers use $t=2$: $a=4$, $b=3$, $c=2$, giving
$4\,\mathrm{Fe}+3\,\mathrm{O_2}\to2\,\mathrm{Fe_2O_3}$. Check: Fe $4=4$, O $6=6$.

\item From $R_2$, $I_1=10-2I_2$. The junction equation $I_1-I_2-I_3=0$ becomes
$10-2I_2-I_2-I_3=0$, i.e. $3I_2+I_3=10$. With $R_3$: $2I_2+3I_3=9$. Solving,
$I_2=3$, $I_3=1$, and $I_1=10-6=4$. Currents $I_1=4$, $I_2=3$, $I_3=1$ amperes. Junction
law: $4=3+1$ \checkmark.

\item Each row operation on the matrix is exactly an equation operation in disguise:
the columns remember which variable each number belongs to, so manipulating rows is the
same as manipulating equations, producing the same solution. The matrix form is
preferred for large systems because it strips away the repeated variable names,
reducing the work to a mechanical, easily automated sequence of arithmetic steps on the
coefficient grid.
\end{solutions}

\end{document}
